\begin{lemma}[Claim 9: zero-divider PE normalization]
Fix the Section 7 minimum-counterexample data and a
\noderef{HierarchyReplacementPresentation} \(P\) on the same exact terminal
tree \(T\).  Assume every proper coloring leaving exactly \(e_0\) uncolored
has a missing color at every vertex.  Suppose the replacement hierarchy has
zero dividers, its terminal series step is PE, and a color \(\alpha\) is
missing on its starred core and at the final vertex \(y_\ell\), but not on
its terminal ladder tree \(T_n\).  Suppose also that no index
\(1\leq i\leq\ell-1\) is already common-missing with \(y_\ell\) under the
coloring of \(P\).

Then either the finite set
\noderef{HierarchyEligibleCommonMissingIndices} is already nonempty, or
there are another replacement presentation \(P'\) and a color \(\beta\)
which is missing under \(P'\) on \(T_n\), on the starred core, and at
\(y_\ell\), while no eligible tail index is common-missing with \(y_\ell\)
under \(P'\).  All ladder trees, hierarchy stages, and terminal-tree edge
identities in \(P'\) are the exact ones in \(P\) and in the originally fixed
counterexample.
\end{lemma}
\begin{proof}
Write \(\mu_n\) for the coloring of \(P\), \(H\) for its hierarchy, and
\(A\) for its last starred stage.  The fields of
\noderef{HierarchyReplacementPresentation} give an actual corresponding
ETT, its maximum property, the full good hierarchy, and the transported
\noderef{HierarchyMinimumCounterexampleContext}.  In particular \(\mu_n\)
is proper, leaves exactly \(e_0\) uncolored, and satisfies the lower-rung,
minimality, and \(2n+11\) hypotheses.

Apply \noderef{HierarchyCounterexamplePenultimateElementary} and
\noderef{HierarchyCounterexampleTailAtLeastTwo} to that transported
context.  Since \noderef{SameGoodHierarchy} fixes the base core, divider
count, stages, and divider positions, its full-tail length and exact prefix
are the originally fixed
\(\ell=\noderef{HierarchyFullTailLength}(C)\) and
\(B=\noderef{HierarchyFullTailPrefix}(C,\ell-1)\).  We obtain
\(2\leq\ell\) and elementarity of \(V(B)\) under \(\mu_n\).

Because \(q=0\) and the terminal kind is PE,
\noderef{HierarchyBase} expresses \(A\) as the exact ordered union
\(T_n\vee R_n\).  The hypotheses place \(\alpha\) in
\(\overline\mu_n(R_n-V(T_n))\).  The good-hierarchy Gamma pairs at index
zero are pairwise disjoint two-element sets indexed by \(D_{n,0}\).
By \noderef{HierarchyStoredColorsCardinality},
\(|D_{n,0}|\leq n\), and hence their union has size at most \(2n\).
The transported large-set inequality gives
\(|\overline\mu_n(T_n)|\geq2n+11\).  Choose
\[
  \beta\in\overline\mu_n(T_n)-\Gamma^0.
\]
This is an actual color; no membership conclusion about either affected
color in \(D_n\) is needed here.

Apply clause (iv) of \noderef{interchange} to \(\alpha\), which is missing
on \(R_n\), and \(\beta\), which is missing on \(T_n\).  It says that their
two-color path components meeting the exact core \(A=T_n\vee R_n\) are
unique.  Properness and missingness make the components through their
first missing vertices path components: at a vertex missing one of the two
colors the corresponding incidence-graph neighbor set has at most one
member.  The component through \(y_\ell\) is therefore disjoint from the
component meeting \(A\); otherwise interchangeability would identify the
two components, while the assumed absence of a common-missing tail index
and elementarity of \(B\) give distinct endpoint components.

Let \(\nu_n\) be the \noderef{TwoColorComponentSwap} on the component
through \(y_\ell\).  We verify, rather than assume, the three transport
conditions of \noderef{StableChangeCasesOneTwoFullTail}.

First, every displayed tail edge carrying \(\alpha\) or \(\beta\) belongs
to the component meeting their first missing vertex in \(A\) or \(B\).
For an ordinary Tashkinov augmentation this follows from the augmentation's
earlier missing-color witness.  The only non-augmentation possibility is
the zero-divider SE/RE first seed isolated by
\noderef{HierarchyTailMissingColorExclusion}; it cannot occur here because
the actual terminal kind is PE.  Hence no exact tail-edge node lies in the
component through \(y_\ell\), and unfolding
\noderef{TwoColorComponentSwap} gives equality of every displayed
tail-edge value.

Second, if the changed component met a vertex of \(B\), properness would
make its maximal alternating path in \(B\) end at a vertex missing one of
\(\alpha,\beta\).  Elementarity of \(B\), together with the first-missing
component in the core, would then identify that endpoint component with
the core component, contradicting disjointness.  Thus the changed component
is outside \(B\).  At every vertex of \(B\), either both incident affected
edges are unchanged or neither is present, so the missing-color set is
unchanged.

Third, colors outside \(\{\alpha,\beta\}\) have identical missing witnesses
by definition of the swap.  The first \(\alpha\)-missing and
\(\beta\)-missing witnesses lie in the unchanged prefix \(B\), so their
least indices are unchanged as well.  This proves equality of
\noderef{HierarchyFirstMissingPosition} for every color in the last
\noderef{HierarchyStoredColorsAt}; it does not assert missing-set equality
at \(y_\ell\).

All remaining hypotheses of \noderef{StableChangeCasesOneTwoFullTail}
are now present after orienting its ordered color pair as
\((\beta,\alpha)\): both colors are missing on \(B\), the selected component
is a \noderef{TwoColorPathComponent}, the component disjointness was proved
above, and its \(q=0\) Cases 1/2 alternative is supplied directly by
\(\beta\in\overline\mu_n(T_n)\).  The component swap is symmetric in the
two named colors, so it is still exactly the coloring \(\nu_n\) constructed
above.  The cited lift returns
an actual corresponding ETT on the full exact tree \(T\), maximum property,
a rebuilt base, a full good hierarchy, and
\noderef{SameGoodHierarchy} data under \(\nu_n\).

Swapping at \(y_\ell\) makes \(\beta\) missing there while leaving its
missing witness on \(T_n\) unchanged.  Thus the rebuilt full tree is still
non-elementary, and these data construct a
\noderef{HierarchyCounterexample}.  Transport the minimum context as
follows.  Rung agreement transfers the lower-rung and fewer-divider
quantifiers.  Same hierarchy stages give equality of divider counts and
of \noderef{HierarchyCounterexampleMeasure}; hence the original
lexicographic-minimality field gives the replacement one.  Correspondence
and terminal stability identify the ladder-tree missing sets, so the
\(2n+11\) inequality transfers.  Properness and exact one-uncoloredness are
fields of the corresponding certificate.  These are precisely all fields
of \noderef{HierarchyMinimumCounterexampleContext}, and therefore give a
new replacement presentation \(P'\).

If \(P'\) has a common-missing tail index, its presentation and color are
immediately witnesses in the filter defining
\noderef{HierarchyEligibleCommonMissingIndices}; this is the first
alternative.  Otherwise the swap has made \(\beta\) missing on \(T_n\),
hence on \(A\), and at \(y_\ell\), and the failed filter condition is exactly
the final universal disjointness assertion.  This is the second alternative.
\end{proof}
